\documentclass[10pt,aspectratio=169]{beamer}
\usepackage[portuguese]{babel}
\usepackage{fontspec,unicode-math}
\setmainfont{Noto Sans}
\setsansfont{Noto Sans}
\setmathfont{STIX Math}
\usepackage{amsmath}
\usepackage{xcolor}
\definecolor{InsperRed}{HTML}{a01924}
\definecolor{Ink}{HTML}{232529}
\setbeamercolor{frametitle}{fg=Ink}
\setbeamercolor{structure}{fg=InsperRed}
\setbeamercolor{normal text}{fg=Ink,bg=white}
\setbeamerfont{frametitle}{size=\large,series=\bfseries}
\setbeamerfont{title}{size=\LARGE,series=\bfseries}
\setbeamertemplate{navigation symbols}{}
\setbeamertemplate{itemize item}{\color{InsperRed}$\bullet$}
\setbeamertemplate{footline}{%
  \hbox{\begin{beamercolorbox}[wd=.84\paperwidth,ht=2.7ex,dp=1.2ex,leftskip=.7cm]{normal text}\scriptsize Inteligência Artificial para Decisão · Insper · 2027-1\end{beamercolorbox}%
  \begin{beamercolorbox}[wd=.16\paperwidth,ht=2.7ex,dp=1.2ex,rightskip=.7cm plus1fil]{normal text}\hfill\scriptsize\insertframenumber\, /\, \inserttotalframenumber\end{beamercolorbox}}}
\setbeamersize{text margin left=12mm,text margin right=12mm}
\begin{document}
\begin{frame}[t]{}
\vspace{14mm}\textcolor{InsperRed}{\bfseries AULA 07}\par\vspace{4mm}{\LARGE\bfseries Retrieval-Augmented Generation}\par\vspace{6mm}{\large Macroeconomia}\par\vspace{12mm}{\color{gray}Paulo C. Marques F. · 2027-1}
\end{frame}

\begin{frame}[t]{Objetivos e verificação}
\small
\begin{itemize}
\item Compreender a arquitetura de RAG e separar erros de ingestão, recuperação e geração, com controle explícito de fontes e datas.
\item Resolver manualmente os exemplos do capítulo, identificando hipóteses e consequências.
\item Especificar a aplicação no Codex gráfico e confrontar ao menos um resultado com evidências.
\end{itemize}
\end{frame}

\begin{frame}[t]{Organização do encontro}
\small
\begin{itemize}
\item Apresentação dos fundamentos e resolução de problemas quantitativos no papel.
\item Laboratório supervisionado e testes produzidos por prompts na interface gráfica do Codex, sem terminal.
\item Atividade individual com dados diferentes e depósito no GitHub antes do quiz seguinte.
\end{itemize}
\end{frame}

\begin{frame}[t]{Recuperação antes da geração}
\small
\begin{itemize}
\item Em Retrieval-Augmented Generation (RAG), a resposta é condicionada a trechos recuperados de um acervo.
\item O mecanismo separa preparação do corpus, indexação, recuperação e geração.
\item Essa distinção importa porque uma resposta pode ser incorreta por duas razões: a evidência não foi recuperada ou foi recuperada, mas o gerador não a utilizou adequadamente.
\item A preparação define documentos, metadados, política de fragmentação (chunking) e possíveis sobreposições.
\end{itemize}
\end{frame}

\begin{frame}[t]{Recuperação antes da geração — continuação}
\small
\begin{itemize}
\item Fragmentos curtos podem perder o referente de expressões como ``essa decisão''; longos demais podem reduzir precisão e ocupar desnecessariamente o contexto.
\item A qualidade não depende de maximizar indiscriminadamente o número de trechos.
\end{itemize}
\end{frame}

\begin{frame}[t]{Recuperação e top-k}
\small
\begin{itemize}
\item Se um índice retorna trechos ordenados por relevância, top-k conserva apenas os primeiros.
\item Considere a ordem e suponha que apenas contém a evidência.
\item Com , o gerador não recebe o trecho correto.
\item Mesmo que escreva uma resposta plausível, ela não estará fundamentada pelo corpus fornecido.
\end{itemize}
\end{frame}

\begin{frame}[t]{Recuperação e top-k — continuação}
\small
\begin{itemize}
\item A recuperação pode ser lexical (termos coincidentes), densa (vetores de embeddings) ou híbrida.
\item A melhor abordagem depende da estrutura dos documentos e das perguntas.
\item Nomes, datas e números exatos são frequentemente úteis para busca lexical; paráfrases podem favorecer busca densa.
\end{itemize}
\end{frame}

\begin{frame}[t]{Citações e distinção entre texto e interpretação}
\small
\begin{itemize}
\item Uma citação documental exige identificador, data ou versão e localização do trecho.
\item Uma resposta sobre determinada reunião do Copom não pode ser sustentada apenas por trecho verdadeiro, porém extraído de reunião diferente.
\item É necessário preservar a dimensão temporal.
\item O sistema deve separar claramente (a) o que o documento afirma, (b) a síntese produzida pela aplicação e (c) interpretações adicionais do analista.
\end{itemize}
\end{frame}

\begin{frame}[t]{Citações e distinção entre texto e interpretação — continuação}
\small
\begin{itemize}
\item Quando o acervo não responde à questão, a abstenção é um resultado legítimo.
\item A saída não deve inventar falas nem atribuir ao Banco Central interpretações não presentes no documento.
\end{itemize}
\end{frame}

\begin{frame}[t]{Segurança documental e avaliação}
\small
\begin{itemize}
\item Um texto recuperado pode conter instruções para o modelo.
\item Essas instruções não têm a autoridade do prompt da aplicação.
\item Deve-se tratar documentos como dados não confiáveis e testar tentativas de prompt injection.
\item Para avaliar o retriever, criar perguntas com trechos relevantes identificados; para avaliar o gerador, conferir cada afirmação com a fonte correta.
\end{itemize}
\end{frame}

\begin{frame}[t]{Segurança documental e avaliação — continuação}
\small
\begin{itemize}
\item No laboratório serão usados documentos didáticos explicitamente fictícios, com registros de reunião e data.
\item Uma extensão opcional substitui esse conjunto por atas oficiais do Banco Central, com proveniência e versões congeladas para reprodução.
\end{itemize}
\end{frame}

\begin{frame}[t]{Exercício para resolver no papel 1}
\small
\small\textbf{7.1 O trecho certo ficou de fora}\par\vspace{4mm}Uma consulta sobre a justificativa do Copom recupera, em ordem, trechos A, B, C, D. Apenas C contém a resposta. Com , quais trechos entram no contexto? O gerador consegue justificar a resposta apenas com o que recebeu? Separe erro de recuperação de erro de geração.
\end{frame}

\begin{frame}[t]{Exercício para resolver no papel 2}
\small
\small\textbf{7.2 Uma referência anacrônica}\par\vspace{4mm}Uma pergunta menciona a decisão de setembro, mas a resposta cita uma ata de julho sem indicar a diferença temporal. O texto citado é verdadeiro em julho. Isso basta para responder à pergunta? Indique os três campos documentais que deveriam acompanhar o trecho.
\end{frame}

\begin{frame}[t]{Laboratório e caso aplicado}
\small
\begin{itemize}
\item Área de aplicação: Macroeconomia.
\item Escolher manualmente a evidência para duas perguntas e registrar reunião/data em cada citação.
\item Solicitar ao Codex extração, fragmentação, indexação local, recuperação top-k e resposta com citações por trecho.
\item Alterar k e tamanho dos fragmentos; observar uma pergunta cuja resposta desaparece do contexto e outra envolvendo reuniões diferentes.
\end{itemize}
\end{frame}

\begin{frame}[t]{Prompt de desenvolvimento — exemplo}
\small
\small\textbf{Instrução fornecida ao Codex:}\par\vspace{5mm}Construa um RAG local sobre as três atas congeladas. Mostre por pergunta os trechos recuperados, reunião, data, localização e resposta fundamentada. Permita abstenção; se houver geração por LLM, use apenas o Codex Python SDK previamente validado; mantenha modo de inspeção de recuperação sem geração.
\end{frame}

\begin{frame}[t]{Critérios de aceitação e validação}
\small
\begin{itemize}
\item Todas as citações apontam para trecho presente na ata correta; respostas sem evidência sinalizam insuficiência; versão de corpus e parâmetros de busca registrados; nenhuma inferência é apresentada como declaração literal do Banco Central.
\item Inspecionar argumentos, saídas e rastros das ferramentas quando houver chamadas externas.
\item Comparar uma saída com a resolução manual do exemplo conceitual.
\end{itemize}
\end{frame}

\begin{frame}[t]{Atividade individual em casa}
\small
\begin{itemize}
\item Duas outras atas oficiais, previamente selecionadas e identificadas, acompanhadas de oito perguntas que exigem distinguir reuniões e justificar citações.
\item Resolver as questões de interpretação, com cálculos e evidência documental.
\item Depositar no GitHub os arquivos solicitados; a entrega habilita o quiz, sem valer nota.
\end{itemize}
\end{frame}

\begin{frame}[t]{Preparação para a avaliação escrita}
\small
\begin{itemize}
\item Exercícios do caderno: 7.1 O trecho certo ficou de fora; 7.2 Uma referência anacrônica; 7.3 Chunking e referência ambígua.
\item Formalizar o raciocínio e interpretar resultados; não será exigida escrita manual de código.
\item No quiz seguinte serão avaliados conceitos desta aula e da atividade feita em casa.
\end{itemize}
\end{frame}

\end{document}
